\documentclass[11pt,a4paper]{article}
\usepackage[T1]{fontenc}
\usepackage[utf8]{inputenc}
\usepackage[margin=2.3cm]{geometry}
\usepackage{booktabs}
\usepackage{longtable}
\usepackage{array}
\newcolumntype{P}[1]{>{\raggedright\arraybackslash}p{#1}}
\usepackage{enumitem}
\usepackage{xcolor}
\usepackage[colorlinks=true,urlcolor=blue!50!black,linkcolor=black]{hyperref}
\setlist{itemsep=2pt,topsep=3pt}
\newcommand{\code}[1]{\texttt{\small #1}}
\sloppy

\title{Moonstone: Audit Pass}
\author{Static review of \code{motleystate/moonstone}, with test-run confirmation}
\date{23 October 2026 (confirmation run added the same day)}

\begin{document}
\maketitle

\section*{Update: confirmation run}

After the static review, the package was installed and the findings were executed. Summary:

\begin{itemize}
  \item \textbf{Every ``to confirm'' finding is confirmed.} Each one raises the error the review predicted
        (Section~2.1 now carries the exact messages).
  \item \textbf{The whole command-line interface (CLI) is dead on a fresh install.} \code{moonstone -{}-help} fails with
        \code{ImportError: cannot import name 'interp' from 'scipy'}, so no option works, including PCA and k-means.
  \item \textbf{New, larger finding: pandas 3.} The declared bound \code{pandas>=2.0.2} lets a fresh install pick up
        pandas 3.1.0, under which 69 of 339 tests fail. With pandas 2.3.3 pinned, 6 tests fail.
  \item \textbf{The Principal Coordinates Analysis (PCoA) tests fail on axis sign only.} Five tests in
        \code{test\_beta.py} expect coordinates whose signs are flipped relative to scikit-bio 0.7.4's output; magnitudes match.
\end{itemize}

\paragraph{Environment.} Linux, CPython 3.12.3, clean virtual environments created with \code{uv}, package installed with
\code{pip install -e .} from \code{master} at commit \code{3930974}. Resolved versions:

\begin{center}
\begin{tabular}{@{}lll@{}}
\toprule
 & \textbf{Environment A (fresh install)} & \textbf{Environment B (pandas pinned)}\\
\midrule
pandas & 3.1.0 & 2.3.3\\
scikit-learn & 1.9.1 & 1.9.1\\
SciPy & 1.15.3 & 1.15.3\\
NumPy & 2.4.6 & 2.4.6\\
scikit-bio & 0.7.4 & 0.7.4\\
\midrule
\code{pytest} result & \textbf{69 failed, 270 passed} & \textbf{6 failed, 333 passed}\\
\bottomrule
\end{tabular}
\end{center}

The classifier findings were exercised by \code{audit/probe\_findings.py}, which calls each library behaviour directly
and then runs \code{ML.svm}, \code{ML.roc\_analysis}, \code{ML.feature\_analysis} and \code{RandomForest.forest} on a
synthetic 40-sample, two-class table. To get past the import error, the probe sets \code{scipy.interp = numpy.interp}
in its own process; no repository code was changed. Results were identical in both environments.

\paragraph{Status at merge (23 October 2026, later the same day).} The findings below describe \code{master} at
\code{3930974}. Since then, pull requests \#113 (prevalence fix) and \#118 (pandas 3 support, and PCoA tests compared up to
axis sign) were merged. On \code{master} at \code{f95b2cb}, the full suite passes in both environments
(340 passed with pandas 3.1.0 and with pandas 2.3.3, optional readers \code{odfpy}, \code{openpyxl} and \code{xlrd}
installed as in CI). Still open: the SciPy import (the CLI still fails at start-up; fix in pull request \#117),
\code{-{}-svm\_classifier} and \code{-{}-random\_forest} (fix-list items 3 and 4), and the classifier rewrite.

\section*{Scope and limits}

This is a first audit pass over the repository \href{https://github.com/motleystate/moonstone}{motleystate/moonstone}
(shallow clone of the default branch, about 16,500 lines of Python, packaged as version 1.1.0).
It aims to answer three questions: what the pipeline looks like, what is likely to work, and where the
machine learning code departs from good practice.

\paragraph{Important limitation (original pass; now lifted, see the update above).} The code was \emph{read, not executed}. Installing the package and running the test
suite was not permitted in this session, because \code{pip install -e .} runs the repository's own
\code{setup.py}. Every finding below comes from reading the source against the dependency versions the project
declares. Findings that depend on library behaviour are marked \emph{to confirm}; the first step of the next session
should be to install in a clean environment and run \code{pytest}, which will confirm or refute them quickly.

\section{What the repository is}

Moonstone is a Python library and command-line interface (CLI) for metagenomics count tables. Its layout:

\begin{itemize}
  \item \code{parsers/}: readers for count tables (Kraken2, MetaPhlAn 2 and 3, QIIME, PICRUSt2) and metadata.
  \item \code{filtering/}, \code{normalization/}: cleaning, filtering and normalization steps that record their history.
  \item \code{analysis/}: diversity (alpha and beta, including Principal Coordinates Analysis, PCoA), statistical tests,
        differential analysis, clustering, and the classifiers (\code{classify.py}, \code{randomForest.py}).
  \item \code{plot/}: bar, box, violin, scatter, heatmap and trendline graphs built on Plotly and Matplotlib.
  \item \code{main.py}: the \code{moonstone} CLI, which dispatches to PCA, k-means, SVM, ROC and random forest.
\end{itemize}

The Kraken2 support is a single parser, \code{SunbeamKraken2Parser}, for the merged table produced by the Sunbeam
pipeline. It does not read native Kraken2 report files directly.

\section{Findings}

\subsection{Code paths that look broken under the declared dependencies}

The project requires pandas $\geq$ 2.0.2 and scikit-learn $\geq$ 1.3.1. Several calls in the classifier modules were
removed in those versions or never existed.

The \textbf{Status} column records the confirmation run (environments A and B gave the same result for every row).

\begin{longtable}{@{}P{3.35cm}P{4.5cm}P{1.7cm}P{5.6cm}@{}}
\toprule
\textbf{Location} & \textbf{Problem} & \textbf{Original confidence} & \textbf{Status and actual error}\\
\midrule
\endhead
\code{randomForest.py:30} & Calls \code{classify.SVM(...)} and then \code{.merge(...)}. The only class in \code{classify.py} is \code{ML}, and it has no \code{merge} method. The \code{-{}-random\_forest} option cannot work. & Certain (grep) & \textbf{Confirmed.} \code{AttributeError: module 'moonstone.analysis.classify' has no attribute 'SVM'}\\
\code{randomForest.py:58} & \code{max\_features="auto"} was removed from \code{RandomForestClassifier} in scikit-learn 1.3. & To confirm & \textbf{Confirmed} (raised at \code{fit}, line 86, once line 30 is bypassed). \code{InvalidParameterError: The 'max\_features' parameter of RandomForestClassifier must be \ldots{} a str among \{'log2', 'sqrt'\} or None. Got 'auto' instead.}\\
\code{classify.py:5} & \code{from scipy import interp} relies on an alias that recent SciPy releases removed. Because \code{main.py} imports \code{classify} at module level, this would break the \emph{whole} CLI at start-up, not only the ROC command. & To confirm (check SciPy version) & \textbf{Confirmed} on SciPy 1.15.3, the newest version the package allows. \code{ImportError: cannot import name 'interp' from 'scipy'}. \code{moonstone -{}-help} fails with this error.\\
\code{classify.py:160} & \code{rfecv.grid\_scores\_} was removed in scikit-learn 1.2 and replaced by \code{cv\_results\_}. & To confirm & \textbf{Confirmed} (library call; not reached end-to-end because line 136 fails first). \code{AttributeError: 'RFECV' object has no attribute 'grid\_scores\_'}\\
\code{classify.py:136} & \code{DataFrame.drop([variable], 1)} passes the axis positionally; pandas 2.0 removed that. & To confirm & \textbf{Confirmed}, pandas 2.3.3 and 3.1.0. \code{TypeError: DataFrame.drop() takes from 1 to 2 positional arguments but 3 were given}. This is the first error \code{-{}-svm\_classifier} hits.\\
\code{classify.py:154} & \code{ascending=[bool]} passes the Python type \code{bool}, not a boolean value. & Certain & \textbf{Confirmed} (library call). \code{ValueError: For argument "ascending" expected type bool, received type type.}\\
\bottomrule
\end{longtable}

\noindent\code{ML.svm} and \code{ML.roc\_analysis} ran to completion once the import was shimmed; their problems are the
methodological ones in Section~2.3, not crashes.

\subsection{New findings from the test run}

\begin{longtable}{@{}P{2.8cm}P{12.5cm}@{}}
\toprule
\textbf{Area} & \textbf{Finding (all observed in the run)}\\
\midrule
\endhead
pandas 3 compatibility & \code{pyproject.toml} allows pandas 3, which a fresh install selects. 63 extra test failures follow, from
removed APIs: \code{DataFrame.applymap} (28 failures; used in \code{parsers/handleCounts.py:46},
\code{normalization/counts/geometric\_mean.py:53}, \code{utils/taxonomy.py:86}), \code{pd.value\_counts} (5;
\code{plot/metadata.py:109,161}, \code{utils/pandas/series.py:179}), \code{fillna(method=...)} (\code{utils/taxonomy.py:44}),
\code{stack(dropna=...)} (3), and the new default string dtype (several assertion mismatches).
\code{handleCounts.py:46} is on the CLI's count-file path, so under pandas 3 the CLI would still fail after the SciPy fix.\\
PCoA tests & 5 failures in \code{tests/analysis/diversity/test\_beta.py} (both environments). Expected and actual coordinates
differ only in sign, e.g.\ expected $(-0.0643, -0.0269)$, got $(0.0643, 0.0269)$. Axis sign in PCoA is arbitrary, so the
tests are brittle rather than the method wrong; this matters for the comparison against the R reference (step 4).\\
Taxa bar-graph test & 1 failure in \code{tests/plot/test\_counts.py} (both environments): the expected first taxon is
\textit{Lactobacillus}, the actual is \textit{Streptococcus}. Likely a tie in ordering; cause \emph{not yet investigated}.\\
\bottomrule
\end{longtable}

\subsection{Code that runs but gives misleading results}

\begin{enumerate}
  \item \textbf{Wrong model evaluated.} In \code{randomForest.py} lines 111--112, the block labelled ``Gradient Boost''
        fits and scores \code{ab} (AdaBoost) instead of \code{gb}. The printed train/test accuracy is AdaBoost's, a second time.
  \item \textbf{Preprocessing leakage.} Scaling is fitted on the whole dataset before cross-validation or splitting
        (\code{classify.py:36--37, 57, 127}; \code{randomForest.py:41}). The held-out folds therefore influence the scaling
        parameters. Scaling should sit inside a scikit-learn \code{Pipeline} so it is refitted on each training fold.
  \item \textbf{Optimistic feature-selection estimate.} Recursive Feature Elimination with Cross-Validation (RFECV) and the final
        Support Vector Machine (SVM) coefficients are computed on all samples, and the cross-validation score reported comes from
        the same data that chose the features. Honest performance needs an outer cross-validation loop around the selection.
  \item \textbf{No handling of class imbalance.} Accuracy is the only metric, there is no \code{class\_weight}, and
        \code{train\_test\_split} in \code{randomForest.py:64} is not stratified. With unbalanced cohorts, accuracy can look high
        for a classifier that mostly predicts the majority class. Balanced accuracy, the Area Under the Receiver Operating
        Characteristic Curve (ROC AUC), and precision/recall per class would be more informative.
  \item \textbf{Not reproducible.} The ROC folds (\code{classify.py:75}), the train/test split, the decision tree, and the
        \code{cross\_val\_score} calls are unseeded. The seeded splitter \code{sf} defined at \code{randomForest.py:63} is not
        used by \code{cross\_val\_score} at line 91.
  \item \textbf{Possible sample dependence (risk).} No grouped splitting is used. If several samples come from the same subject,
        site or sequencing batch, ordinary k-fold lets near-duplicates sit in both training and test folds. Whether this applies
        depends on the cohort; \code{GroupKFold} would remove the risk.
  \item \textbf{Biased importances.} Random forest importances are impurity-based and computed on the 75\% training split only.
        Permutation importance on held-out data is more trustworthy for ranking taxa.
\end{enumerate}

\subsection{PCoA biplot}

The biplot is built in \code{analysis/diversity/beta.py} from scikit-bio's \code{pcoa} and \code{pcoa\_biplot}. The helper
\code{\_scale\_biplot} rescales and shifts each axis separately so that feature arrows fit the range of the sample cloud.
This is a reasonable visual overlay, but it changes the angles between arrows and does not place them at the origin, so the
arrows are not true biplot vectors. The method should be documented as an overlay, or replaced by a single common scaling.
(The function also edits its input DataFrame in place; this is harmless today because the caller only uses the index afterwards,
but it is fragile.)

\subsection{Tests and continuous integration}

Twenty source modules have no matching test file, including exactly the ones with the defects above:
\code{classify.py}, \code{randomForest.py}, \code{clustering.py}, \code{main.py}, the MetaPhlAn parser, and the box, violin,
heatmap and histogram graphs. That explains how broken calls survived. Continuous integration runs \code{pytest} on Python 3.9
and 3.13 and only fails the lint step on syntax errors and undefined names, so a missing attribute such as \code{classify.SVM}
passes unnoticed.

\subsection{Packaging}

Dependencies are declared twice, in \code{pyproject.toml} and \code{setup.py}, with different bounds (for example numpy
$\geq$ 1.26.4 versus $\geq$ 1.24.3, and SciPy with and without an upper limit). \code{requirements.txt} is a
\code{pip-compile} output that includes macOS-only and obsolete packages (\code{appnope}, \code{backcall}). One source of
truth would remove this ambiguity.

\section{Ranked fix list (after the confirmation run)}

Ranked by how much each fix unblocks, then by effort. Costs are \emph{rough} estimates of agent usage at Claude Opus 5.5
list prices (\$4 per million input tokens, \$20 per million output, \$0.20 per million cached input); they are not measured,
the mapping from these prices to your cloud credit is not known to me, and real cost can easily be twice the estimate if a
fix needs several review rounds.

\begin{longtable}{@{}P{0.4cm}P{9.6cm}P{2.2cm}P{2.5cm}@{}}
\toprule
\textbf{\#} & \textbf{Fix} & \textbf{Size} & \textbf{Rough cost}\\
\midrule
\endhead
1 & Replace \code{scipy.interp} with \code{numpy.interp} (\code{classify.py:5,88}); add a CLI smoke test that runs
\code{moonstone -{}-help}. Restores the whole CLI under pandas 2. & One line plus one test & \$1--3\\
2 & pandas 3: fix the call sites listed in Section~2.2 and the string-dtype test expectations, \emph{or} cap
\code{pandas<3} as a stop-gap. Fixing is better; capping changes a dependency bound and should be your call. Add pandas 3 to
the continuous integration (CI) matrix. & Moderate & \$5--12 (fix) or \$1--2 (cap)\\
3 & \code{-{}-svm\_classifier}: \code{drop(..., axis=1)} at line 136, \code{ascending=True} at line 154,
\code{cv\_results\_["mean\_test\_score"]} at line 160; one test per method. & Small & \$3--6\\
4 & \code{-{}-random\_forest}: use \code{MergeDF} in \code{get\_matrix}, \code{max\_features="sqrt"} (the old
\code{"auto"} meant \code{"sqrt"} for classifiers, so results stay comparable), fit and score \code{gb} not \code{ab} at lines
111--112; one test. & Small & \$3--6\\
5 & PCoA tests: compare coordinates up to sign per axis, so they pass on any scikit-bio version. Natural to fold into step 4
of the project (R reference). & Small & \$2--4\\
6 & Single source of truth for dependencies (drop \code{install\_requires} from \code{setup.py}); this is what let pandas 3 in
unnoticed. & Small & \$1--3\\
7 & Classifier rewrite (project step 3): pipeline, stratified/grouped seeded splits, class weights, balanced accuracy, nested
cross-validation, permutation importance, tests. & Large & \$25--60\\
\bottomrule
\end{longtable}

Items 1--6 together: roughly \$15--35. With item 7: roughly \$40--95. This confirmation step itself was a mid-sized session
on the same scale as items 2--4 combined.

\section{Suggested next steps, in order (original pass)}

\begin{enumerate}
  \item In a clean Python 3.11--3.13 environment, install the package and run \code{pytest}; record which ``to confirm'' items
        above are real.
  \item Fix the import-time and attribute errors (Section 2.1) and add one smoke test per CLI option so they cannot regress.
  \item Rewrite the classifiers around one shared, tested module: \code{Pipeline} with scaling inside, stratified and
        seeded splits, optional grouping, class weights, balanced accuracy and ROC AUC, nested cross-validation for feature
        selection, and permutation importance.
  \item Add tests for the PCoA biplot against the R reference output you already have.
  \item Write a short methods page for biologists: what each metric means, why the splits are done this way, and what a
        trustworthy result looks like.
\end{enumerate}

\end{document}
